\chapter{Introduction} \label{chap:one}
\section{Research background} \label{sec:1.1}
Near-surface soils play a critical role in the performance of many geotechnical infrastructure systems. The hydromechanical state of such soils is influenced by both atmospheric conditions and applied stresses. Earthworks, pavement foundations, railway embankments, cuttings, landfill covers, and shallow containment systems are commonly constructed or reworked under unsaturated conditions. Many railway assets in the United Kingdom were constructed long before modern compaction control, drainage design and unsaturated soil mechanics were available as routine engineering tools. Their long-term serviceability depends on the ability of the soil mass to tolerate cycles of rainfall infiltration, evaporation, vegetation-driven water uptake, shrinkage, swelling and traffic loading. Field evidence from clay slopes has shown that seasonal soil water content and pore-water pressure variations can be related quantitatively to rainfall, evapotranspiration and vegetation effects, and that such cycles can contribute to shrink-swell movements, strain softening and progressive instability (Smethurst et al., 2012). For formation layers of railways and roads, the same environmental cycles interact with repeated traffic loads and can alter both water retention response and cyclic mechanical behaviour (Azizi et al., 2023).



Consequently, their performance depends not only on the density, grading, and stress state established during construction, but also on the subsequent migration of water through the pore network, seasonal fluctuations in suction, and the progressive microstructural changes induced by repeated drying–wetting cycles.



\section{Research objectives} \label{sec:1.2}
In response to the challenges outlined above, the overall goal of this thesis is to advance xxxxx. To achieve this goal, this thesis investigates key aspects of xxxxxx, with specific objectives listed as follows:
\begin{itemize}
  \item xxxx
  \item xxxx
  \item xxxx
  \item xxxx
\end{itemize}

\section{Thesis overview}
This thesis consists of seven chapters. The remaining chapters are outlined below. \par